\chapter{Dérivation, variations et optimisation}\label{C06}
\begin{repere}{Le parcours du chapitre}
\textbf{Objectifs.} Calculer une dérivée sur son domaine ; étudier variations et extrema ; comparer des courbes ; optimiser ; exploiter les symétries.

\textbf{Prérequis.} Nombre dérivé, calcul algébrique, signe d’un trinôme, intervalles.

\textbf{Repère annuel.} 1–10 décembre : 8 h, consolidation les 15 et 17\end{repere}
\section*{Diagnostic : six questions, 10 minutes}\label{C06-D}
Répondre sans consulter le cours. Indiquer les questions pour lesquelles une aide est nécessaire. Ce diagnostic ne donne pas lieu à une note.
\begin{enumerate}
\item Dérivée de $x^2$ en 3 ?
\item Domaine de $\dfrac{1}{x-2}$ ?
\item Signe de $(x-1)(x+2)$ ?
\item Développer $(x^2+1)(x-2)$.
\item Pour $f(x)=x^2$, comparer $f(-x)$ et $f(x)$.
\item Si $f(1)=2$, $f\prime(1)=0$, équation de la tangente ?
\end{enumerate}
\paragraph{Après le diagnostic.} Reprendre les notions signalées avec le professeur ; refaire une question analogue, puis expliquer la règle utilisée.
\paragraph{Oral et logique.} Présenter un tableau de variations en distinguant signe de f et signe de sa dérivée. Expliquer pourquoi les bornes d’un problème d’optimisation doivent être examinées.
\clearpage
\section*{Chercher avant de formaliser}\label{C06-RECH}
Une fonction croissante peut-elle avoir une tangente horizontale ? Chercher un exemple ; peut-elle avoir une dérivée strictement négative en un point intérieur ?
\paragraph{Organisation.} Recherche individuelle 5 min, confrontation à deux 5 min, mise en commun 5 min. Garder une trace d’un essai, même non abouti. Une conjecture doit ensuite être justifiée.
\section{Fonction dérivée et fonctions usuelles}\label{C06-U1}
\begin{repere}{Définition et résultats}
Une fonction est dérivable sur un intervalle lorsqu’elle l’est en chacun de ses points ; sa dérivée associe à x le nombre $f\prime(x)$. Les dérivées usuelles sont : constante $\mapsto0$ ; $x\mapsto1$ ; $x^2\mapsto2x$ ; $x^3\mapsto3x^2$ ; $\dfrac{1}{x}\mapsto -\dfrac{1}{x^2}$ pour $x\ne0$ ; $\sqrt x\mapsto \dfrac{1}{2\sqrt x}$ pour $x>0$. Pour $n\in\Z$, $(x^n)\prime=nx^{n-1}$ sur son domaine de dérivabilité ; pour $n=0$, la dérivée est 0 directement.
\end{repere}
\paragraph{Statut et justification.} Carré : le taux vaut $2a+h$. Inverse : pour $a\ne0$, $h\ne0$ et $a+h\ne0$, $\dfrac{[\dfrac{1}{a+h}-\dfrac{1}{a}]}{h}=-\dfrac{1}{[a(a+h)]}$, qui tend vers $-\dfrac{1}{a^2}$. Les autres formules sont admises ici.
\exemple{Pour $f(x)=x^{-3}$, le domaine est $\R\setminus\{0\}$ et $f\prime (x)=-3x^{-4}=-\dfrac{3}{x^4}$. Pour $g(x)=\sqrt x$, la fonction existe en 0, sa dérivée usuelle seulement pour $x>0$.}\label{C06-EX1}
\begin{erreur}Le domaine de dérivabilité peut être plus petit que le domaine de définition.\end{erreur}
\application{Dériver $x^5$, $x^{-2}$ et $\sqrt x$ en précisant les domaines.}\label{C06-AP1}
\section{Somme, produit, inverse et quotient}\label{C06-U2}
\begin{repere}{Définition et résultats}
Pour u et v dérivables, $(u+v)\prime=u\prime+v\prime$, $(\lambda u)\prime=\lambda u\prime$, $(uv)\prime=u\prime v+uv\prime$. Là où $v\ne0$, $(\dfrac{1}{v})\prime =-\dfrac{v\prime}{v^2}$ et $(\dfrac{u}{v})\prime =\dfrac{u\prime v-uv\prime}{v^2}$.
\end{repere}
\paragraph{Statut et justification.} Pour le produit, ajouter et retrancher $u(a+h)v(a)$ au numérateur. Le taux devient $\dfrac{u(a+h)[v(a+h)-v(a)]}{h}+\dfrac{v(a)[u(a+h)-u(a)]}{h}$. Comme une fonction dérivable est continue (admis), on obtient $u(a)v\prime(a)+v(a)u\prime(a)$. Les autres règles sont admises.
\exemple{Pour $f(x)=\dfrac{x^2+1}{x-1}$, $x\ne1$, $f\prime (x)=\dfrac{[2x(x-1)-(x^2+1)]}{(x-1)^2}=\dfrac{x^2-2x-1}{(x-1)^2}$. Garder les parenthèses évite de changer un seul signe au lieu de tout le numérateur.}\label{C06-EX2}
\begin{erreur}La dérivée d’un produit n’est pas le produit des dérivées ; la dérivée d’un quotient n’est pas leur quotient.\end{erreur}
\application{Dériver $(x^2+1)(3x-2)$ puis $\dfrac{1}{x^2+1}$.}\label{C06-AP2}
\section{Signe de la dérivée et extrema}\label{C06-U3}
\begin{repere}{Définition et résultats}
Sur un intervalle I, si $f\prime\geq0$, f est croissante ; si $f\prime\leq0$, elle est décroissante ; si $f\prime=0$ partout, elle est constante. Une dérivée strictement positive entraîne une stricte croissance. Si f possède un extremum local en un point intérieur où elle est dérivable, sa dérivée y est nulle. La réciproque est fausse.
\end{repere}
\paragraph{Statut et justification.} Théorèmes de variation admis. Pour un minimum intérieur en a, les taux à droite sont positifs ou nuls, ceux à gauche négatifs ou nuls ; s’ils convergent vers le même réel, il est 0. Même argument pour un maximum.
\exemple{$f(x)=x^3-3x$ : $f\prime(x)=3(x-1)(x+1)$. f croît sur $]-\infty;-1]$, décroît sur $[-1;1]$, croît sur $[1;+\infty[$ ; maximum local $f(-1)=2$, minimum local $f(1)=-2$. Aucun de ces extrema n’est global sur $\R$.}\label{C06-EX3}
\par\medskip\noindent\begin{minipage}{\linewidth}\centering
\begin{tikzpicture}\begin{axis}[width=.88\linewidth,height=6.4cm,axis lines=middle,grid=major,grid style={gray!25},tick label style={font=\normalsize},label style={font=\normalsize},clip=false,xmin=-2.3,xmax=2.3,ymin=-4,ymax=4,xlabel={$x$},ylabel={$f(x)$}]\addplot[bleu,thick,domain=-2.2:2.2,samples=100]{x^3-3*x};\addplot[only marks,vert] coordinates {(-1,2)(1,-2)};\end{axis}\end{tikzpicture}
\par\medskip $f(x)=x^3-3x$ : un maximum local en $-1$ et un minimum local en 1.\end{minipage}\par\medskip

\begin{erreur}Le théorème s’applique sur un intervalle. Une fonction peut décroître sur deux intervalles séparés sans décroître sur leur réunion.\end{erreur}
\application{Étudier $f(x)=x^2-4x+1$ sur $[0;3]$ et donner ses extrema sur cet intervalle.}\label{C06-AP3}
\section{Parité, inégalités et comparaison}\label{C06-U4}
\begin{repere}{Définition et résultats}
Sur un domaine symétrique par rapport à 0, f est paire si $f(-x)=f(x)$ et impaire si $f(-x)=-f(x)$. La courbe d’une fonction paire est symétrique par rapport à l’axe des ordonnées ; celle d’une impaire par rapport à l’origine. Pour comparer f et g, on étudie le signe de $f-g$, éventuellement avec sa dérivée.
\end{repere}
\paragraph{Statut et justification.} Les symétries découlent des coordonnées des points images. Pour une inégalité, une étude de variations peut donner un minimum de la différence, puis son signe ; la dérivée ne se substitue pas directement à cette différence.
\exemple{Pour $x>0$, $h(x)=x+\dfrac{1}{x}-2$ a pour dérivée $\dfrac{x^2-1}{x^2}$. Elle décroît jusqu’à 1 puis croît ; $h(1)=0$, donc $x+\dfrac{1}{x}\geq 2$. L’égalité est obtenue seulement en 1.}\label{C06-EX4}
\begin{erreur}Verifier le domaine avant la parité. Le signe de $f\prime-g\prime$ renseigne sur la variation de $f-g$, pas directement sur son signe.\end{erreur}
\application{Étudier la parité de $x^4+x^2$ et $x^3-x$ ; démontrer $x^2+1\geq2x$.}\label{C06-AP4}
\section{Optimisation et méthode de Newton}\label{C06-U5}
\begin{repere}{Définition et résultats}
Optimiser consiste à définir la variable et son domaine, exprimer la grandeur à optimiser, puis comparer toutes les valeurs candidates, y compris les bornes. La méthode de Newton utilise la tangente pour approcher une solution de $f(x)=0$ : $x_{n+1}=x_n-\dfrac{f(x_n)}{f}\prime (x_n)$ lorsque $f\prime(x_n)\ne0$. Elle n’est utilisée que dans des cas favorables.
\end{repere}
\paragraph{Statut et justification.} La formule de Newton vient de l’intersection avec l’axe horizontal de $y=f(x_n)+f\prime(x_n)(x-x_n)$. La convergence n’est pas garantie en général ; on la constate et contrôle dans le cas retenu.
\exemple{Pour $f(x)=x^2-2$ avec $x_0=1$, Newton donne $x_1=\dfrac{3}{2}$, $x_2=\dfrac{17}{12}$. Les carrés sont 2,25 puis environ 2,00694 : on approche $\sqrt2$. Le programme doit contrôler dérivée nulle et nombre maximal d’étapes.}\label{C06-EX5}
\begin{erreur}Une valeur numérique fournie par un algorithme n’est pas une preuve de minimum ni de convergence.\end{erreur}
\application{Un rectangle de périmètre 12 m a largeur x. Déterminer l’aire maximale par dérivation.}\label{C06-AP5}
\section*{Parcours d’exercices}\label{C06-PARC}
Commencer par les exercices 1 et 2. Choisir ensuite un exercice 3 ou 4 sans aide. Les exercices 5 à 7 sont une réserve de transfert, de recherche ou de programmation : le professeur en désigne un. L’exercice 8 se reprend à distance.
\begingroup
\exercice{C06-E1}{Une collection de dérivées}
On considère les fonctions\[f(x)=3x^4-2x+7,\quad g(x)=(x^2+2)(x-1),\quad h(x)=\dfrac{2x+1}{x+3}.\]
\begin{enumerate}[label=\arabic*.]
\item Préciser le domaine de définition de chacune de ces fonctions.
\item Calculer leurs dérivées en nommant, à chaque fois, la règle utilisée. Pour $g$, conserver d’abord les deux termes de la règle du produit ; pour $h$, écrire le numérateur avant de le réduire.
\item Contrôler le résultat concernant $g$ en développant son expression avant de dériver.
\end{enumerate}
\exercice{C06-E2}{Un tableau complet}
La fonction $f$ est définie sur $[0;4]$ par $f(x)=x^3-6x^2+9x+1$. On cherche ses valeurs minimale et maximale sur cet intervalle.
\begin{enumerate}[label=\arabic*.]
\item Calculer et factoriser $f\prime(x)$, puis étudier son signe sur $[0;4]$.
\item Dresser le tableau de variations de $f$ en indiquant les valeurs aux bornes et aux points où la dérivée s’annule.
\item Donner le minimum et le maximum de $f$ sur l’intervalle, ainsi que toutes les abscisses où ils sont atteints.
\end{enumerate}
\exercice{C06-E3}{Un quotient}
Pour tout réel strictement positif $x$, on pose $f(x)=x+\dfrac4x$.
\begin{enumerate}[label=\arabic*.]
\item Étudier les variations de $f$ sur $]0;+\infty[$. Présenter les facteurs qui déterminent le signe de la dérivée.
\item En déduire la plus petite valeur de $x+\dfrac4x$, puis écrire une inégalité valable pour tout $x>0$. Préciser exactement son cas d’égalité.
\end{enumerate}
\exercice{C06-E4}{Comparer deux courbes}
Les fonctions $f(x)=x^3$ et $g(x)=x$ sont définies sur $\R$. On note $\mathcal C_f$ et $\mathcal C_g$ leurs courbes dans un même repère.
\begin{enumerate}[label=\arabic*.]
\item Déterminer les points d’intersection de ces deux courbes.
\item Étudier leur position relative sur $\R$ en choisissant une méthode adaptée. Traduire votre résultat par les intervalles où $\mathcal C_f$ est au-dessus ou au-dessous de $\mathcal C_g$.
\end{enumerate}
\exercice{C06-E5}{Un emballage}
On fabrique une boîte sans couvercle à partir d’une feuille carrée de 12 cm de côté. On découpe dans chacun des quatre coins un carré de côté $x$ cm, puis on relève les bords restants. L’épaisseur de la feuille est négligée.
\begin{enumerate}[label=\arabic*.]
\item Dessiner le patron et exprimer la hauteur ainsi que les deux dimensions de la base de la boîte. Préciser les valeurs de $x$ qui donnent une boîte de volume non nul.
\item Exprimer le volume $V(x)$ en centimètres cubes. Déterminer les dimensions qui rendent ce volume maximal.
\item Contrôler les valeurs aux deux bornes du domaine et donner le volume maximal exact.
\end{enumerate}
\par\medskip\noindent\begin{minipage}{\linewidth}\centering
\begin{tikzpicture}[x=.5cm,y=.5cm,>=Stealth]\draw[thick](0,0)rectangle(12,12);\foreach \x/\y in {0/0,10/0,0/10,10/10}{\fill[ocre!20](\x,\y)rectangle++(2,2);\draw[ocre,thick](\x,\y)rectangle++(2,2);}\draw[dashed](2,2)rectangle(10,10);\draw[<->](0,13)--node[above]{12 cm}(12,13);\draw[<->](-1,0)--node[left]{$x$}(-1,2);\draw[<->](0,-1)--node[below]{$x$}(2,-1);\node at(6,6){base};\end{tikzpicture}
\par\medskip Les quatre zones colorées sont retirées ; les pointillés indiquent les plis. Le dessin représente une valeur illustrative de x et ne donne pas l’optimum.\end{minipage}\par\medskip

\exercice{C06-E6}{Stationnaire sans extremum}
Un élève affirme : « Si la dérivée s’annule en un point, la fonction possède forcément un extremum local en ce point. » On étudie la fonction $f(x)=x^3$ pour tester cette affirmation.
\begin{enumerate}[label=\arabic*.]
\item Calculer $f\prime(0)$.
\item Comparer à $f(0)$ les valeurs de $f(x)$ pour des réels $x$ négatifs puis positifs, aussi proches de 0 que l’on veut.
\item Conclure sur l’existence d’un extremum local en 0 et sur l’affirmation proposée.
\end{enumerate}
\exercice{C06-E7}{Comparer des stratégies}
On souhaite démontrer que, pour tout réel $x$, $x^2-6x+12\geq3$.
\begin{enumerate}[label=\arabic*.]
\item Proposer une démonstration en faisant apparaître un carré.
\item Proposer une seconde démonstration en étudiant les variations d’une fonction.
\item Préciser le cas d’égalité. Comparer les deux méthodes : laquelle demande ici le moins de calculs ?
\end{enumerate}
\exercice{C06-E8}{Reprise différée}
La fonction inverse $f(x)=\dfrac1x$ a pour domaine $D=\R\setminus\{0\}$. Un élève calcule $f\prime(x)=-\dfrac1{x^2}<0$ et conclut que $f$ est décroissante sur tout $D$.
\begin{enumerate}[label=\arabic*.]
\item Identifier l’hypothèse manquante dans l’application du théorème reliant signe de la dérivée et variations.
\item Indiquer les deux intervalles sur lesquels sa conclusion est correcte.
\item À l’aide des nombres $-1$ et 1, construire un contre-exemple à sa conclusion sur tout $D$.
\end{enumerate}
\paragraph{Travail personnel.} Effectuer uniquement la tâche désignée à l’issue de la séance, dans la progression datée : 10–20 min maximum. Les exercices dépendent des notions déjà étudiées ; certaines questions sont donc différées. Reprendre l’exercice 8 une semaine après le chapitre (10 min). Si le temps prévu est dépassé, noter la difficulté et apporter les essais ; ne pas accumuler toute la réserve.
\endgroup
\Needspace{12\baselineskip}\section*{À faire dans le manuel CQFD}
Ces exercices se traitent dans votre manuel. La progression indique la séance ; ils remplacent le travail personnel prévu ce jour-là (10 à 20 minutes au total).
\paragraph{CQFD, exercice 12 p. 156, questions 1 et 2.}\label{C06-M1} Calculer la différence, étudier son signe et traduire celui-ci en position relative des courbes.
\paragraph{CQFD, exercice 16 p. 156, questions f et g.}\label{C06-M2} Donner pour chaque fonction un intervalle où elle est à la fois négative et croissante ; justifier les deux propriétés séparément.
\begingroup\small
\clearpage\section*{QCM 1 : comprendre et choisir}\label{C06-Q1}
Cinq questions pédagogiques à réponse unique. Choisir une réponse par question ; expliquer ensuite un distracteur. Ce format de chapitre ne reproduit pas le format officiel de l’EAM.
\begin{enumerate}
\item $(x^3)\prime$ vaut
\begin{enumerate}[label=\Alph*.,itemsep=2pt]
\item $x^2$
\item $3x^2$
\item $3x$
\item $\dfrac{x^4}{4}$
\end{enumerate}
\item $(\dfrac{1}{x})\prime $ vaut
\begin{enumerate}[label=\Alph*.,itemsep=2pt]
\item $-\dfrac{1}{x^2}$
\item $\dfrac{1}{x^2}$
\item $-\dfrac{1}{x}$
\item $1$
\end{enumerate}
\item $(uv)\prime$ vaut
\begin{enumerate}[label=\Alph*.,itemsep=2pt]
\item $u\prime v\prime$
\item $u\prime+v\prime$
\item $uv\prime$
\item $u\prime v+uv\prime$
\end{enumerate}
\item $f\prime>0$ sur un intervalle implique
\begin{enumerate}[label=\Alph*.,itemsep=2pt]
\item f positive
\item f négative
\item f strictement croissante
\item f constante
\end{enumerate}
\item $\dfrac{x^2+1}{x-1}$ est définie si
\begin{enumerate}[label=\Alph*.,itemsep=2pt]
\item $x>0$
\item $x\ne1$
\item $x\ne-1$
\item $x\ne0$
\end{enumerate}
\end{enumerate}
\endgroup
\begingroup\small
\Needspace{18\baselineskip}\section*{QCM 2 : comprendre et choisir}\label{C06-Q2}
Cinq questions pédagogiques à réponse unique. Choisir une réponse par question ; expliquer ensuite un distracteur. Ce format de chapitre ne reproduit pas le format officiel de l’EAM.
\begin{enumerate}
\item $x^3-x$ est
\begin{enumerate}[label=\Alph*.,itemsep=2pt]
\item paire
\item constante
\item impaire
\item ni paire ni impaire
\end{enumerate}
\item $f\prime(a)=0$ implique nécessairement
\begin{enumerate}[label=\Alph*.,itemsep=2pt]
\item tangente horizontale
\item maximum
\item minimum
\item f constante
\end{enumerate}
\item Dérivée de $\sqrt x$ :
\begin{enumerate}[label=\Alph*.,itemsep=2pt]
\item $2\sqrt x$
\item $\dfrac{1}{\sqrt x}$
\item $\dfrac{\sqrt x}{2}$
\item $\dfrac{1}{2\sqrt x}$, $x>0$
\end{enumerate}
\item Minimum de $x+\dfrac{1}{x}$ sur $]0;+\infty[$ :
\begin{enumerate}[label=\Alph*.,itemsep=2pt]
\item $0$
\item $2$
\item $1$
\item aucun
\end{enumerate}
\item Newton pour $x^2-2$ n’est pas défini avec $x_0=$
\begin{enumerate}[label=\Alph*.,itemsep=2pt]
\item $1$
\item $2$
\item $0$
\item $-1$
\end{enumerate}
\end{enumerate}
\subsection*{Vrai ou faux ? Justifier.}\label{C06-VF}
\begin{enumerate}
\item La dérivée d’une fonction positive est positive.
\item Un extremum intérieur d’une fonction dérivable donne une dérivée nulle.
\item Le produit de deux fonctions croissantes est croissant.
\end{enumerate}
\begin{repere}{Bilan personnel}Je dérive sur le bon domaine ; je simplifie sans perdre les valeurs interdites ; je distingue extremum local et global ; je choisis parfois une factorisation plutôt qu’une dérivée.

Pour chaque objectif, indiquer : acquis sans aide ; réussi avec aide ; à retravailler. Écrire un exemple et une erreur à éviter.\end{repere}
\endgroup
\ifprof\corrsection{6}\label{C06-CORR}
\textbf{Usage professeur.} Les réponses qui suivent concernent les tâches non résolues dans le cours. Faire d’abord expliciter la démarche ; fournir une aide graduée avant la solution. Les durées sont des estimations à étalonner en classe.
\subsection*{Diagnostic et remédiation}
\begin{enumerate}
\item $6$. \textit{Action :} Reprendre le taux $2a+h$.
\item $\R\setminus\{2\}$. \textit{Action :} Annuler le dénominateur.
\item Positif à l’extérieur de $[-2;1]$, négatif entre les racines, nul aux bornes. \textit{Action :} Construire deux lignes de signes.
\item $x^3-2x^2+x-2$. \textit{Action :} Distribuer chaque terme sans omission.
\item Ils sont égaux. \textit{Action :} Substituer -x entre parenthèses.
\item $y=2$. \textit{Action :} Reprendre l’équation tangente et simplifier.
\end{enumerate}
\subsection*{Activité de recherche}
$x^3$ est strictement croissante avec tangente horizontale en 0. Pour une fonction croissante dérivable, les taux de part et d’autre sont positifs ou nuls, donc la dérivée ne peut être négative.
\subsection*{Applications immédiates}
\begin{enumerate}
\item $5x^4$ sur $\R$ ; $-\dfrac{2}{x^3}$ pour $x\ne0$ ; $\dfrac{1}{2\sqrt x}$ pour $x>0$.
\item $2x(3x-2)+3(x^2+1)=9x^2-4x+3$ ; $-\dfrac{2x}{(x^2+1)^2}$, sur $\R$.
\item $f\prime=2x-4$, donc décroissance jusqu’à 2 puis croissance. Minimum $-3$ en 2 ; valeurs aux bornes 1 et $-2$, donc maximum 1 en 0.
\item Première paire, seconde impaire ; différence $x^2+1-2x=(x-1)^2\geq0$, égalité en 1.
\item $A(x)=x(6-x)$ sur $[0;6]$ ; $A\prime=6-2x$. Maximum 9 m$^2$ en 3.
\end{enumerate}
\subsection*{Exercices : démarche et résultats}
\paragraph{Exercice 6.1 -- Une collection de dérivées.} $12x^3-2$ ; $2x(x-1)+(x^2+2)=3x^2-2x+2$ ; $\dfrac{[2(x+3)-(2x+1)]}{(x+3)^2}=\dfrac{5}{(x+3)^2}$ pour $x\ne-3$.
\paragraph{Exercice 6.2 -- Un tableau complet.} $f\prime=3(x-1)(x-3)$ ; signes $+,-,+$. Valeurs $f(0)=1$, $f(1)=5$, $f(3)=1$, $f(4)=5$. Minimum 1 en 0 et 3 ; maximum 5 en 1 et 4. Les bornes comptent.
\paragraph{Exercice 6.3 -- Un quotient.} $f\prime =1-\dfrac{4}{x^2}=\dfrac{(x-2)(x+2)}{x^2}$. Signe négatif avant 2, positif après. Minimum 4 en 2 ; pour $x>0$, $x+\dfrac{4}{x}\geq 4$, égalité seulement en 2.
\paragraph{Exercice 6.4 -- Comparer deux courbes.} $f-g=x(x-1)(x+1)$. Signes : négatif sur $]-\infty;-1[$, positif sur $]-1;0[$, négatif sur $]0;1[$, positif sur $]1;+\infty[$. Intersection en $-1,0,1$. On peut faire le tableau des trois facteurs sans dérivation.
\paragraph{Exercice 6.5 -- Un emballage.} $V=x(12-2x)^2=4x(6-x)^2$. $V\prime=12(x-2)(x-6)$, positif sur $[0;2[$, négatif sur $]2;6[$. $V(0)=V(6)=0$, $V(2)=128$ cm$^3$. Maximum pour x=2 cm ; boîte $8\times8\times2$.
\paragraph{Exercice 6.6 -- Stationnaire sans extremum.} $f\prime=3x^2$, donc valeur 0 en 0. Dans tout voisinage de 0, des x négatifs donnent $f(x)<0$ et des x positifs $f(x)>0$ ; aucun extremum.
\paragraph{Exercice 6.7 -- Comparer des stratégies.} Carré : $(x-3)^2+3\geq3$. Dérivée : $2x-6$, minimum en 3, valeur 3. La première évite une étude de variation complète ; la seconde illustre le lien entre chapitres.
\paragraph{Exercice 6.8 -- Reprise différée.} Elle décroît sur $]-\infty;0[$ et $]0;+\infty[$. Elle ne décroît pas sur leur réunion : $-1<1$ mais $f(-1)=-1<1=f(1)$. Le domaine n’est pas un intervalle.
\paragraph{CQFD, exercice 12 p. 156, questions 1 et 2.} $f(x)-g(x)=\dfrac12x^2-1=\dfrac12(x-\sqrt2)(x+\sqrt2)$. La courbe de f est strictement au-dessus de celle de g pour $x<-\sqrt2$ ou $x>\sqrt2$, au-dessous entre ces deux nombres ; égalité aux abscisses $\pm\sqrt2$.
\paragraph{CQFD, exercice 16 p. 156, questions f et g.} Pour $f(x)=x^3$, on peut choisir $]-\infty;0[$ : $f(x)<0$ et la fonction est strictement croissante. Pour $g(x)=x^2-4$, on peut choisir $[0;2[$ : $g(x)<0$ et la fonction est strictement croissante sur cet intervalle. Des sous-intervalles non réduits à un point conviennent aussi.
\subsection*{QCM : réponses et raisons}
\textbf{QCM 1.}
\begin{enumerate}
\item B. Puissance entière.
\item A. Taux de l’inverse.
\item D. Règle du produit.
\item C. Signe de dérivée et variation.
\item B. Dénominateur non nul.
\end{enumerate}
\textbf{QCM 2.}
\begin{enumerate}
\item C. $f(-x)=-f(x)$.
\item A. Les autres conclusions demandent plus.
\item D. Formule et domaine.
\item B. Atteint en 1.
\item C. Dérivée $2x_0=0$.
\end{enumerate}
\subsection*{Vrai ou faux}
\begin{enumerate}
\item Faux : $f(x)=\dfrac{1}{x}>0$ sur $]0;+\infty[$, mais $f\prime(x)<0$.
\item Vrai, avec les hypothèses précisées ; aux bornes, la conclusion n’est pas imposée.
\item Faux sans hypothèses de signe : sur $[-2;-1]$, $u(x)=v(x)=x$ croissent mais $x^2$ décroît.
\end{enumerate}
\paragraph{Différenciation.} Pour un blocage technique, proposer la première étape seulement (domaine, facteur ou univers de référence), puis retirer l’aide. Pour un élève autonome, demander une seconde méthode et le traitement d’un cas limite. À l’oral, autoriser une reformulation initiale dans la langue maîtrisée, puis exiger une phrase mathématique précise en français.
\fi
